% Purpose: Demonstrate a fluid, self-contained seismic-field dictionary.
% Status: Draft template; schema and human layout review pending.
% Source-of-truth: .claude/agents/latex-docxer.md (extracted schema archetype).
% Origin: Illustrative field assumptions; no repository schema was inspected.
\documentclass{article}
\usepackage[a4paper,margin=2cm]{geometry}
\usepackage{amsmath}
\usepackage{booktabs,tabularx,threeparttable}

\begin{document}
\section*{Illustrative schema dictionary}

\begin{table}[htbp]
  \centering
  \small
  \begin{threeparttable}
    \caption{Illustrative seismic-field dictionary.
      The rows distinguish types, domains, and optional fields.
      All constraints are schematic assumptions, not verified project
      schema definitions; replace them using an inspected schema anchor.}
    \label{tab:data-dictionary}
    \begin{tabularx}{\linewidth}{@{} l l l c X @{}}
      \toprule
      \textbf{Field} & \textbf{Type} & \textbf{Domain}
        & \textbf{Null?\tnote{a}} & \textbf{Meaning / Binding} \\
      \midrule
      \texttt{network} & String & Network code & No
        & Network identifier; verify code format. \\
      \texttt{station} & String & Station code & No
        & Station identifier; verify code format. \\
      \texttt{channel} & String & Channel code & No
        & Channel identifier; verify code format. \\
      \texttt{phase} & Enum & $\{P,S\}$ & No
        & Illustrative compressional/shear phase labels. \\
      \texttt{onset\_time} & DateTime & ISO 8601 UTC & No
        & Arrival time; identify observed, analyst, or predicted origin. \\
      \texttt{probability} & Float & $[0,1]$ & Yes
        & Model score; verify meaning and calibration.\tnote{b} \\
      \texttt{residual} & Float & Signed seconds & Yes
        & Travel-time residual; verify sign convention. \\
      \texttt{weight} & Integer & $\{0,1,2,3,4\}$ & Yes
        & Illustrative analyst quality code; verify its domain. \\
      \bottomrule
    \end{tabularx}
    \begin{tablenotes}[flushleft]
      \footnotesize
      \item[a] Nullability is illustrative, not a rule for any pipeline stage.
      ISO 8601 denotes a date/time representation; UTC is Coordinated
      Universal Time. Residuals use seconds.
      \item[b] A score is a model prediction, not an analyst label or
      established probability calibration.
      \item Source/schema/version/reviewer: [TODO: provide authorized anchors].
    \end{tablenotes}
  \end{threeparttable}
\end{table}
\end{document}
